
In this section, we present the formalization of adaptivity for a given program through three steps using the new methodology proposed in Section~\ref{sec:intro-methodology-adapt-formalization}.
Each step corresponds to one of the three procedures of the flow-chart in the ``Adaptivity Formalization'' block at the middle of Figure~\ref{fig:adaptfun-structure}. We
summarize them as follows.
 \begin{enumerate}
 \item The first step in 
 Section~\ref{sec:dynamic-datadep} defines the reachability sub-property of 
 the \emph{adaptivity} as a
 variable \emph{may-dependency} relation based on the trace semantics.
% in Section~\ref{sec:language-os}.
% This is the reachability perspective of this property.
 %
 \item In Section~\ref{sec:dynamic-graph}, 
 we define the quantitative sub-property for a program's \emph{adaptivity} as the weight of each labeled variable.
 Then we construct a weighted dependency graph namely
the \emph{semantics-based dependency graph} of a program $c$, which carries both 
sub-properties. 

\item 
% The last step is the intuitive \emph{adaptivity} quantity formalization presented in Section~\ref{sec:dynamic-adapt}.
 Based on the \emph{semantics-based dependency graph},
 we define the formal \emph{adaptivity} in definition~\ref{def:trace_adapt} as a quantitative reachability hyper-property in Section~\ref{sec:dynamic-adapt}.
\end{enumerate}
% , this is the second major part of this adaptivity analysis framework
% built on the language design. 
Through two examples in Section~\ref{sec:dynamic-examples}, I show that the formalized
adaptivity matches the program's intuitive \emph{adaptivity}.
%  more precisely and efficiently than previous works.

